% ============================================================
%  02_advocacy_brief.tex
%  Autonomy, Capacity, and the Right to Exit
%  ADVOCACY BRIEF --- For disability rights orgs, civil liberties groups,
%  patient advocacy organizations
%  Justin Bogner, 2026
%  Compile: xelatex 02_advocacy_brief.tex (twice)
% ============================================================

\documentclass[11pt, letterpaper]{article}
\usepackage{campaign}
\usepackage{xspace}
\usepackage{colortbl}

\campaignheader%
  {Autonomy, Capacity, and the Right to Exit}%
  {Advocacy Brief --- Bogner, 2026}

\geometry{
  letterpaper,
  left=1.15in, right=1.15in,
  top=1.1in, bottom=1.2in,
  headheight=22pt
}

\titleformat{\section}
  {\sffamily\large\bfseries\color{C@accent}}
  {}{0em}{}[{\color{C@rulegray}\titlerule[0.5pt]}]

\titleformat{\subsection}
  {\sffamily\normalsize\bfseries\color{C@darkgray}}
  {}{0em}{}

\begin{document}

\campaigntitle%
  {AUTONOMY, CAPACITY, AND THE RIGHT TO EXIT}%
  {An Advocacy Brief on Capacity-Based Medical Aid in Dying}%
  {Advocacy Brief}%
  {Justin Bogner}%
  {2026}

\suitenote

\begin{policybox}[title={\sffamily\bfseries\small\color{C@accent}
  Who This Brief Is For}]
This brief is written for disability rights organizations, civil liberties advocacy groups,
patient rights advocates, and policy organizations that work at the intersection of
autonomy and social provision. It takes the disability rights critique of expanded MAiD
as a serious argument deserving serious engagement---not dismissal---and offers a
structural response grounded in the same values that animate that critique.
\end{policybox}

\bigskip

\section{Starting Point: What the Critics Are Right About}

The disability rights organizations opposing expanded MAiD eligibility---DAWN Canada, ADAPT,
the UN Special Rapporteur on the Rights of Persons with Disabilities---have identified a
real and documentable problem. It deserves to be stated plainly before the response begins.

Health Canada's 2024 Annual Report documents 732 Track~2 provisions (non-foreseeable death
cases), representing a 17\% year-over-year increase. Of these, 61.5\% of recipients reported
disability as a condition---nearly double the 31.6\% rate in terminal Track~1 cases. The
Ontario Chief Coroner's reviews have documented individual cases in which inadequate housing,
income insufficiency, and absence of community care appear as primary conditions motivating
MAiD requests.

These are not fringe cases. They are a pattern. And the disability rights organizations are
correct to name what that pattern suggests: in some instances, the Canadian state is
processing deaths that are attributable not to irremediable medical conditions but to
irremediable policy failures.

\begin{findingbox}[title={\sffamily\bfseries\small\color{C@accent}
  The Concern, Stated Fairly}]
Expanding MAiD eligibility to non-terminal cases does not extend autonomy if the conditions
under which people make that choice are coercive by virtue of deprivation. A person who
chooses death because they cannot access adequate housing, disability support, or social
connection is not exercising free choice in any meaningful sense. They are being failed by
the state and given a bureaucratically clean mechanism to exit that failure.

This is a legitimate concern. It demands a structural response, not a procedural footnote.
\end{findingbox}

\section{Where the Protectionist Response Fails}

The policy position that follows from the disability rights concern is typically:
\textit{restrict eligibility, close the exit, compel continued existence until adequate
supports are in place.}

This position has three fatal problems.

\subsection{It does not resolve the underlying conditions}

The Protectionist Model---restrictions on MAiD eligibility as the mechanism for compelling
social investment---has governed most jurisdictions for decades. Canada's existing Track~2
framework, limited as it is, has been operative since 2021. The United States has restricted
MAiD to terminal cases since the first state laws in the 1990s. The United Kingdom has no
MAiD provision at all.

In those same decades: disability support funding has been chronically underfunded across
all these jurisdictions. Social housing stocks have declined. Community mental health
infrastructure has been cut. Institutional waiting lists have grown. The Protectionist
position did not produce the dignified conditions its defenders claim it would compel.
There is no historical evidence that locking the exit creates political pressure for better
lives. The evidence is that it creates administrative invisibility for people whose suffering
is politically inconvenient to see.

\subsection{It substitutes one coercion for another}

When the state denies a person with intact decisional capacity the right to act on their
informed preference---on the grounds that the state disapproves of the conditions producing
that preference---it is not protecting that person. It is overriding their expressed will
with institutional judgment. This is precisely the form of paternalism that disability rights
advocacy has historically opposed: the substitution of external authority for personal agency
in the lives of people the state deems vulnerable.

The person who cannot access adequate housing and finds their life intolerable is not made
better off by being prohibited from acting on that assessment. They are worse off, in a way
that is now legally compelled. They remain in the conditions that produced the crisis, with
the additional indignity of the state having decided it knows better than they do whether
those conditions are sufficient to justify exit.

\subsection{It makes the state's failure invisible}

This is the most consequential problem. Under current restrictive frameworks, when a person
in inadequate conditions is denied MAiD access and continues living in those conditions,
that outcome is recorded as nothing. There is no document. No report. No attributed failure.
The person's continued existence, however inadequate, is counted as a policy success.

When that same person eventually dies---in a hospital, in inadequate housing, by
unassisted means---the causal chain between policy failure and death is severed by the
absence of documentation. The state is never held accountable, because the mechanism for
accountability does not exist.

\section{The Forcing Function: A Structural Answer}

The response this framework offers to the disability rights concern is not a procedural
footnote. It is a structural mechanism designed to make state failure politically legible.

\begin{findingbox}[title={\sffamily\bfseries\small\color{C@accent}
  The Core Mechanism}]
Under this framework, where an applicant identifies inadequate housing, income, disability
support, or social connection as a primary or contributing condition of their MAiD
preference, the administering authority is \textbf{required} to transmit a formal written
report to the relevant government ministry documenting these conditions as a driver of the
provision. This reporting obligation is mandatory, not discretionary, and must be fulfilled
within 30 days of provision.
\end{findingbox}

The policy rationale is straightforward.

\subsection{State failure becomes documented public record}

Aggregated across a population, mandatory socioeconomic reporting produces a dataset that
did not previously exist: the number of people, by year, by region, by condition, who
received MAiD in circumstances attributable in part to policy failure. That dataset is
a political instrument. It is citable, attributable, and specific.

The disability rights organizations have been arguing, correctly, that inadequate supports
are driving some MAiD uptake. This framework makes that argument documentable. It transforms
an advocacy claim into a government data series.

\subsection{Political cost creates structural incentive}

A government that must publish, annually, the number of its citizens who died with inadequate
housing as a documented driver faces a political cost that the current system does not
impose. The current system produces no such document. Inadequate deaths are absorbed into
aggregate mortality statistics and disappear.

The Right to Exit, combined with mandatory socioeconomic documentation, functions as a
consumer protection standard for existence: the state must either provide conditions worth
living in or produce a public record of why it did not. This is the incentive structure
for social investment that the Protectionist position cannot create, because it requires
the deaths to be visible.

\subsection{The alternative to coercion is accountability, not access}

The disability rights organizations want the state to improve conditions before unlocking
the exit. The framework presented here wants both simultaneously: unlock the exit, and
document every choice made in the shadow of state failure. These are not in conflict.
They are complementary. The exit cannot remain locked as the condition for investment
in life, because that bet has already been lost. The door has been locked. The investment
has not arrived.

\section{Robust Safeguards for Vulnerable Populations}

Acknowledging that some MAiD uptake is driven by socioeconomic deprivation does not mean
accepting it uncritically. The procedural architecture of this framework includes specific
protections for exactly the populations the disability rights community is concerned about.

\subsection{Mandatory alternatives presentation}

Evaluators are required to present, in documented specificity, every available alternative:
housing assistance programs, income support, disability services, community connection,
peer support, palliative care, psychiatric intervention. Not a pamphlet. A documented
conversation with referrals offered and declinations recorded. An applicant who does not
know that housing assistance exists cannot give an informed refusal of housing assistance.
The evaluation must close that gap.

\subsection{Independent coercion screening}

Every assessment includes a confidential interview with a social worker or patient advocate
who has no institutional stake in the outcome. This person's role is to look for external
pressure---from family, from financial circumstances, from the care system---and to provide
the applicant with space to express concerns they may not feel safe expressing in the
presence of their treating team.

\subsection{Absolute revocation}

The right to revoke is unconditional and available until the final moment. No provision
date may be locked in. No bureaucratic or social momentum may constrain it. This protects
against the ``inertia coercion'' that can occur when a process has begun and withdrawal
feels socially costly.

\subsection{Socioeconomic documentation is mandatory, not optional}

The reporting obligation is not a clinician's discretionary note. It is a statutory
requirement. A provision that occurs without this report, where the conditions warranted
it, constitutes a failure of the administering authority. This is enforceable.

\section{The Values We Share}

This framework and the disability rights critique it engages are animated by the same
underlying values. Both hold that:

\begin{itemize}
  \item The state has an obligation to provide conditions in which disabled people,
    socioeconomically marginal people, and isolated people can live with dignity.
  \item Choices made under conditions of deprivation are not fully free choices, and
    the state bears responsibility for those conditions.
  \item Disabled people are not less capable of autonomous decision-making than
    non-disabled people, and should not be subjected to heightened paternalistic
    oversight on that basis.
  \item The substitution of institutional judgment for personal agency in the lives
    of people deemed vulnerable is a form of harm, not protection.
\end{itemize}

The disagreement is not about values. It is about which mechanism best advances those
values given the evidence about what restrictive frameworks have actually produced.

\subsection{The point of genuine disagreement}

The disability rights position holds that the risk of the state normalizing death as a
response to deprivation is severe enough to justify categorical restriction, even at the
cost of overriding the expressed preferences of people with genuine capacity. This framework
holds that categorical restriction, given the historical evidence, does not produce the
investment it is supposed to compel, and that overriding expressed autonomous preferences
on a population-wide basis is itself a serious harm that must be weighed.

This is a genuine disagreement about empirical predictions and about how to weigh competing
harms. It deserves honest engagement, not mutual dismissal.

\section{A Proposed Basis for Dialogue}

The following points represent the minimal shared ground from which a productive policy
conversation between this framework and disability rights organizations might begin:

\begin{enumerate}
  \item \textbf{Agree on the documentation obligation.} Even organizations opposed to
    expanded eligibility should be able to support mandatory socioeconomic reporting as
    an improvement on the status quo. If deaths are occurring due to inadequate support
    under any eligibility framework, those deaths should be documented. This is not
    controversial in principle.

  \item \textbf{Agree on the alternatives mandate.} Mandatory documented presentation
    of all available alternatives, with referrals offered and declinations recorded,
    advances the goals of disability rights advocacy regardless of eligibility framework.
    It ensures that no one proceeds to MAiD without genuine knowledge of what was available.

  \item \textbf{Engage on the empirical question.} Has the restrictive framework produced
    the social investment it was intended to compel? If not, on what basis should the
    restriction be maintained? This question deserves honest engagement with the historical
    record, not assumption.

  \item \textbf{Disagree on eligibility while cooperating on process.} Even where
    eligibility remains contested, process improvements---independent coercion screening,
    absolute revocation, documentation requirements---can be pursued cooperatively.
\end{enumerate}

\begin{policybox}[title={\sffamily\bfseries\small\color{C@accent}
  A Direct Statement to Disability Rights Organizations}]
You have been right that the state is failing people and that some MAiD uptake reflects
that failure. This framework does not dispute that. It disputes only the conclusion that
the remedy is restriction rather than accountability.

If your position is that disabled people with intact decisional capacity should not have
the same right to self-determination that non-disabled people have---because their choices
might be influenced by conditions the state has failed to provide---then you are arguing
for a form of protective paternalism that disability rights advocacy has historically and
correctly opposed in every other domain.

The door must stand unlocked. What must change is the system of accountability that operates
on the other side of it. This framework builds that accountability system. We are, on that
goal, allies.
\end{policybox}

\bigskip

{\small\sffamily\color{C@midgray}
\textbf{Key Sources:}
Health Canada, \textit{Sixth Annual Report on MAiD in Canada} (2024 data; released November 2025);
DAWN Canada, \textit{MAiD and the Inadequacy of Social Supports} (2023);
Ontario Office of the Chief Coroner, \textit{MAiD Case Reviews} (2022--2025);
UN Special Rapporteur on the Rights of Persons with Disabilities (2019);
\textit{Carter v.\ Canada (AG)}, 2015 SCC 5.}

\end{document}
